\documentclass[12pt]{article}
\usepackage[utf8]{inputenc}
\usepackage[T1]{fontenc}
\usepackage[margin=1in]{geometry}
\usepackage{setspace}
\usepackage{xcolor}
\usepackage[colorlinks=true, urlcolor=blue!60!black, linkcolor=blue!60!black]{hyperref}
\setlength{\parindent}{2em}
\setlength{\parskip}{1.25em}
\renewcommand{\baselinestretch}{1.0}
\begin{document}
\begin{center}
{\LARGE \textit{NLRB v. Noel Canning}, 573 U.S. 513 (2014)\par}
\vspace{0.8em}
{\large Case Brief --- Four-Page Edition\par}
\end{center}
\vspace{1em}

\section{Memory Jogger}

The Recess Appointments Clause reaches intra-session as well as inter-session recesses and vacancies that arose before the recess, but a recess of three days or less is too short and one of fewer than 10 days is presumptively too short. The Senate is in session when it says it is, so President Obama's January 2012 NLRB appointments, made during a three-day gap between pro forma sessions, were invalid (9--0 in judgment) (\href{https://www.law.cornell.edu/supremecourt/text/12-1281}{Cornell LII}).

\section{Detailed Case Facts}

The NLRB found that Noel Canning, a Pepsi-Cola bottler in Yakima, Washington, unlawfully refused to execute a collective-bargaining agreement. Noel Canning petitioned the D.C. Circuit, arguing that the Board lacked the three lawfully appointed members required for a quorum under \textit{New Process Steel} because Sharon Block, Richard Griffin, and Terence Flynn had been invalidly appointed. President Obama had nominated all three in 2011. On December 17, 2011, the Senate resolved to hold ``pro forma session[s],'' with ``no business . . . transacted,'' every Tuesday and Friday until January 23, 2012. On January 4, 2012, during the three-day break between the January 3 and January 6 sessions, the President recess-appointed the three.

\section{Procedural History}

The D.C. Circuit vacated the Board's order, holding that ``the Recess of the Senate'' means only inter-session recesses and that ``Vacancies that may happen during the Recess'' means only vacancies first arising then. The Supreme Court granted certiorari, added the question of pro forma sessions, and affirmed on that ground on June 26, 2014.

\section{Judicial Votes}

Affirmed, \textbf{9--0} in judgment; \textbf{5--4} on the Clause's scope.
\begin{itemize}
\item \textbf{Majority:} Justice Breyer, joined by Justices Kennedy, Ginsburg, Sotomayor, and Kagan.
\item \textbf{Concurring in the judgment:} Justice Scalia, joined by Chief Justice Roberts and Justices Thomas and Alito.
\item \textbf{Dissents:} None.
\end{itemize}

\section{Holding}

(1) The Clause applies to both inter-session and intra-session recesses, but a break of three days or fewer is too short and one under 10 days is presumptively too short absent very unusual circumstances. (2) The President may fill any vacancy that exists during a qualifying recess, even if it arose earlier. (3) The Senate is in session when it says it is, if under its own rules it retains the capacity to transact business. Because the Senate was in session on January 3 and 6, the three-day break was too short, and the appointments were unconstitutional.

\section{Analysis of Opinions}

\textbf{Majority (Breyer): method.} The Clause is a ``subsidiary'' mechanism to keep the executive functioning while the Senate is unavailable. Invoking Madison's idea of ``liquidation,'' \textit{McCulloch}, \textit{The Pocket Veto Case}, and Frankfurter's \textit{Youngstown} concurrence, Breyer gave ``great weight'' to long historical practice where the text is ambiguous.

\textbf{Intra-session recesses.} Dictionaries and usage of the founding era applied ``the recess'' to intra-session breaks, and the Senate is equally unavailable during them. Presidents since Andrew Johnson have made thousands of such appointments without formal Senate objection. But the Adjournments Clause treats three days as \emph{de minimis}, and no appointment in two centuries came during a break under 10 days. Only extraordinary emergencies, not Senate political opposition, could overcome that presumption.

\textbf{Pre-existing vacancies.} Jefferson read ``happen'' as ``happen to be,'' and Attorney General Wirt as ``happen to exist.'' The narrow reading would paralyze offices, and practice since Madison and the 1940 Pay Act amendments confirm the broad one.

\textbf{Pro forma sessions.} Declining to pierce the Senate record with C-SPAN footage, the Court followed the Rules of Proceedings Clause, \textit{United States v. Ballin}, and \textit{Marshall Field v. Clark}. The Senate could act by unanimous consent, and it actually passed a payroll-tax extension during a December 23 pro forma session.

\textbf{Scalia (concurring in the judgment).} Separation-of-powers rules protect liberty and cannot be changed by acquiescence (\textit{Chadha}, \textit{Free Enterprise Fund}). ``The Recess,'' paired with ``the next Session,'' means only the inter-session break, and ``happen'' means arise. The 10-day presumption is judicial legislation, and an ``adverse possession'' theory favors presidential aggrandizement.

\section{Examples: Future Applications}

\begin{enumerate}
\item \textbf{Same-Side (Invalid): Pro Forma August.} An appointment during a three-day gap between pro forma sessions is invalid, and an eight-day break is presumptively too short; Senate obstruction is no ``unusual circumstance.''
\item \textbf{Same-Side (Void Order): Lost Quorum.} A commission member recess-appointed during a five-day break never lawfully held office, so the commission's order, issued without a quorum, must be set aside.
\item \textbf{Opposite-Side (Valid): An 18-Day Summer Recess.} After both Houses adjourn for 18 days without pro forma sessions, the President may fill a Cabinet vacancy that arose earlier.
\item \textbf{Opposite-Side (Valid): Between Sessions.} An appointment during a 16-day inter-session break to fill an eight-month-old ambassadorial vacancy is valid, though Scalia's reading would doom it.
\item \textbf{Fence-Sitter: Catastrophe in a Nine-Day Recess.} A terrorist attack kills the Secretary of Homeland Security and her deputy. The Court left room for a national catastrophe but has never applied the exception.
\end{enumerate}

\section{Critique}

\textbf{A methodological paradox.} Breyer is functionalist on the first two questions, rejecting formal distinctions between kinds of recesses, yet formalist on the third, accepting nominal pro forma sessions at face value.

\textbf{Adverse possession and bright lines.} Scalia's critique remains potent: because a unitary President acts decisively while a multimember Senate faces collective-action problems, treating silence as acquiescence ratchets toward executive power. And neither ``3 days'' nor ``10 days'' appears in Article II, \S~2, cl. 3, illustrating the tension between administrability and textual fidelity.

\section{Key Quotations}

\begin{quote}
``We hold that, for purposes of the Recess Appointments Clause, the Senate is in session when it says it is, provided that, under its own rules, it retains the capacity to transact Senate business.''
\end{quote}
\noindent --- Justice Breyer, for the Court, 573 U.S. at 550.

\begin{quote}
``In sum, we conclude that the phrase `the recess' applies to both intra-session and inter-session recesses. If a Senate recess is so short that it does not require the consent of the House, it is too short to trigger the Recess Appointments Clause. See Art. I, \S~5, cl. 4. And a recess lasting less than 10 days is presumptively too short as well.''
\end{quote}
\noindent --- Justice Breyer, for the Court, 573 U.S. at 538.

\begin{quote}
``We have not previously interpreted the Clause, and, when doing so for the first time in more than 200 years, we must hesitate to upset the compromises and working arrangements that the elected branches of Government themselves have reached.''
\end{quote}
\noindent --- Justice Breyer, for the Court, 573 U.S. at 526.

\end{document}
